\chapter{군과 대칭}

\begin{learninggoal}
군의 정의를 이해하고, 대칭과 가역적 변환을 군으로 모델링한다.
\end{learninggoal}

\section{왜 군인가}
군은 한 번 적용한 뒤 되돌릴 수 있는 변환들의 합성 구조를 추상화한다. 이 장은 순열, 정수의 덧셈, 합동류, 행렬군을 출발점으로 군의 공리와 기본 예제를 전개한다.

\section{집필 예정}
\begin{itemize}
  \item 이항연산과 군의 공리
  \item 부분군과 생성부분군
  \item 순환군과 원소의 위수
  \item 잉여류와 라그랑주 정리
  \item 정규부분군과 몫군
  \item 군 준동형과 동형정리
  \item 군 작용과 궤도--안정자 정리
\end{itemize}
